
Um Spiele, Lehrerinnen und Lehrer, sowie die Items in das Programm einzubinden, haben wir diese als Objekte implementiert da sie so gut mit den in der Lösungsidee beschriebenen Eigenschaften versehen werden konnten. Dafür haben wir das PalyerObject, das LehrerObjekt und das ItemType gecodet.\\
\subsection{PlayerObject}
Das PalyerObject erweitert das BaseObjekt und übernimmt so wichtige Eigenschaften. Im PlayerObject haben wir die Klasse PlayerClass erstellt, um Platz für die Funktionen der Gestaltung der Figur, die Funktionen des Kampfes und das Inventar zu schaffen. Das Inventar ist eine Reihung, welche als neues, in der Länge unbegrenztes Array initialisiert wurde. Am beim gewinnen eines Minispiels wird, wie oben erwähnt, ein Item dem Array hinzugefügt.\\

\begin{lstlisting}package czg.objects;

import czg.scenes.BaseScene;
import czg.scenes.InventarScene;
import czg.scenes.SceneStack;
import czg.util.Capsule;
import czg.util.Draw;
import czg.util.Images;
import czg.util.character_creator.SaveFile;

import java.awt.*;
import java.awt.image.BufferedImage;
import java.util.*;
import java.util.List;
import java.util.function.Function;

import static czg.MainWindow.PIXEL_SCALE;

/**
 * Die Spielfigur
 */
public class PlayerObject extends BaseObject{

    //Anlegen einer Reihung "Inventar", in welchem die Items, auf die der Spieler zugreifen kann, gespeichert werden
    public final List<ItemObject> inventar = new ArrayList<>();

\end{lstlisting}

Danach folgt der Code, welcher es dem Anwender oder der Anwenderin des Programmes ermöglicht, die Spielerfigur zu gestalten und die Spielfiguren in die Szene zeichnet, welcher jedoch von Ashley erstellt wurde und in der Bedienungsanleitung erklärt wird.\\

\begin{lstlisting}
    

    // Standardfarben
    public static final SaveFile defaultColors = new SaveFile(
            new Color(57, 37, 35),  // Haare: Braun
            new Color(241, 241, 44),// Haut: Gelb
            new Color(45, 96, 145), // Hoodie: Blau
            new Color(97, 105, 112) // Hose: Grau
    );

    /**
     * Farbe der Haare
     */
    public final Capsule<Color> haare = new Capsule<>(defaultColors.haare());
    /**
     * Farbe des Hoodies
     */
    public final Capsule<Color> hoodie = new Capsule<>(defaultColors.hoodie());
    /**
     * Farbe der Hose
     */
    public final Capsule<Color> hose = new Capsule<>(defaultColors.hose());
    /**
     * Farbe der Haut
     */
    public final Capsule<Color> haut = new Capsule<>(defaultColors.haut());

    /**
     * Zuordnung [Farbe im Original-Sprite] -> [Speicherort der Farbe, durch die sie ersetzt werden soll]
     */
    private static final Map<Integer, Function<PlayerObject,Capsule<Color>>> spriteColorMap = new HashMap<>();
    static {
        spriteColorMap.put(0xFF75FB4C, p->p.haare);
        spriteColorMap.put(0xFFEA33F7, p->p.hoodie);
        spriteColorMap.put(0xFF0000F5, p->p.hose);
        spriteColorMap.put(0xFFEA3323, p->p.haut);
    }

    /**
     * Singleton-Instanz
     */
    public static final PlayerObject INSTANCE = new PlayerObject();


    /**
     * Privater Konstruktor, der nur für {@link #INSTANCE} verwendet wird
     */
    private PlayerObject() {
        // Vorläufig ohne Bild initialisieren
        super(null);

        // Vorlage laden und Größe bestimmen
        Image baseSprite = Images.get("/assets/characters/PlayerBase.png");
        width = baseSprite.getWidth(null) * PIXEL_SCALE;
        height = baseSprite.getHeight(null) * PIXEL_SCALE;

        // Tatsächlichen Sprite als zunächst leeres Bild erstellen
        sprite = new BufferedImage(width, height, BufferedImage.TYPE_INT_ARGB);
        // Farben anwenden
        updateSprite();
    }
    
    //Funktion zum Festlegen einer zufälligen x-Koordinate für die Spieler-Figur
    public static int GetRandomX(){
        Random r = new Random();
        int min = 35;
        int max = 790;
        return r.nextInt((max - min) + 1) + min;
    }


\end{lstlisting}

\begin{lstlisting}
    
  
    /**
     * Farben {@link #haare}, {@link #haut}, {@link #hoodie} und {@link #hose} anwenden
     */
    public void updateSprite() {
        // Vorlage laden
        BufferedImage template = (BufferedImage) Images.get("/assets/characters/PlayerBase.png");
        // Grafik-Objekt zum Zeichnen erstellen
        Graphics2D g = (Graphics2D) sprite.getGraphics();

        // Alles löschen
        g.setColor(Draw.TRANSPARENCY);
        g.fillRect(0,0,template.getWidth(null), template.getHeight(null));

        for(int x = 0; x < template.getWidth(null); x++) {
            for(int y = 0; y < template.getHeight(null); y++) {
                // Farbe aus der Vorlage abfragen
                int templateColor = template.getRGB(x,y);

                // Die eingestellte Farbe auswählen ...
                if(spriteColorMap.containsKey(templateColor))
                    g.setColor(spriteColorMap.get(templateColor).apply(this).get());
                // ... oder, wenn es sich nicht um ein Pixel handelt,
                // welches nicht umgefärbt werden soll, die Farbe aus
                // der Vorlage beibehalten
                else
                    g.setColor(new Color(templateColor, true));

                // Skaliertes Pixel zeichnen
                g.fillRect(x*PIXEL_SCALE,y*PIXEL_SCALE,PIXEL_SCALE,PIXEL_SCALE);
            }
        }
    }


\end{lstlisting}

Nach der von Emil erstellten Methode für den Kampf folgt die @Override Fuktion, welche pro Bildaktualisierung mit einer if-Funktion abfragt, ob auf die Spielerfigur geklickt wurde. Falls ja, wird die Inventarscene mit der im bereits definierten push-Funktion nach oben geschoben, das Inventar also geöffnet.
\\
\begin{lstlisting}
    
    @Override
    public void update(BaseScene scene) {
        // Inventar Öffnen, wenn die Figur angeklickt wird
        if(isClicked())
            SceneStack.INSTANCE.push(new InventarScene());
    }
    
}

\end{lstlisting}

\subsubsection{InventarScene}
Die Inventarszene erwitert die BaseScene und übernimmt so wieder ihre Eigenschaften. Die InventarScene hat im gegensatz zu den anderen Szenen keine mit ihr verknüpfte Bilddatei, sondern wird beim öffnen immer wieder neu gezeichnet. Die größe, Position auf dem Bildschirm und der Abstand zwischen den aufgeführten Items werden durch statische Integer initialisiert, da sich diese nicht verändern sollen.\\

\begin{lstlisting}
    package czg.scenes;

import czg.objects.BaseObject;
import czg.objects.ButtonObject;
import czg.objects.ItemObject;
import czg.objects.PlayerObject;
import czg.util.Draw;
import czg.util.Images;

import java.awt.*;

import static czg.MainWindow.*;

public class InventarScene extends BaseScene {

    /**
     * X-Koordinate der linken Seite
     */
    private static final int iLeft = (int)(WIDTH * 0.03);
    /**
     * Y-Koordinate der oberen Seite
     */
    private static final int iTop = (int)(HEIGHT * 0.03);
    /**
     * Breite des Inventars
     */
    private static final int iWidth = WIDTH - 2 * iLeft;
    /**
     * Höhe des Inventars
     */
    private static final int iHeight = (int)(HEIGHT * 0.3);
    /**
     * Abstand zwischen Items
     */
    private static final int iPadding = (int)(WIDTH * 0.025);


\end{lstlisting}
Daraufhin wird der Constructor InventarScene initiiert, in welchen das ButtonObject "exit" lädt. Dieser Button schließt mit der pop-Funktion das Inventar.
\\
\begin{lstlisting}
      public InventarScene() {
        // Button zum Schließen
        ButtonObject exit = new ButtonObject (
                Images.get("/assets/minigames/general/button_exit.png"),
                SceneStack.INSTANCE::pop
        );

        exit.x = iLeft + iWidth - iPadding - exit.sprite.getWidth(null) * PIXEL_SCALE;
        exit.y = iTop + (iHeight / 2) - (exit.sprite.getHeight(null) * PIXEL_SCALE )/2;
        objects.add(exit);
    }
\end{lstlisting}

In der @Override-Funktion wird das eigentliche Inventar was zu sehen ist mit der draw-Funktion gezeichnet. Hierfür wird eine Farbe und eine Größe für das Rechteck festgelegt, welches die Basis bildet und eine Farbe und eine Schriftart für den Text, welcher die Namen der aufgelisteten Items anzeigt. Danach folgt eine Schleife, welche jedes Element des Inventar-Arrays der Spielerklasse erfasst und so mit Namen und Bild in im Padding erfassten Abstand unter Berücksichtigung der Bildgröße zeichnet. Um die Items richtig zu zeichnen, werden die Bilddateien mit dem Faktor PixelScale multipliziert.
\\
\begin{lstlisting}
     @Override
    public void draw(Graphics2D g) {
        // Hintergrund
        g.setColor(new Color(0x9A6B9C));
        g.fillRect(iLeft, iTop, iWidth, iHeight);

        // Für Text
        g.setColor(Color.WHITE);
        g.setFont(Draw.FONT_TITLE.deriveFont(20f));

        // Items mit Beschriftung zeichnen
        int x = iLeft + iPadding;
        for (int i = 0; i < PlayerObject.INSTANCE.inventar.size(); i++) {
            // Item abfragen
            ItemType item = PlayerObject.INSTANCE.inventar.get(i);

            // Sprite zeichnen
            int y = iTop + ((int) (HEIGHT * 0.05));
            g.drawImage(item.SPRITE, x, y, item.SPRITE.getWidth(null) * PIXEL_SCALE, item.SPRITE.getHeight(null) * PIXEL_SCALE, null);

            // Text zeichnen
            Draw.drawTextCentered(g, item.NAME, x + item.SPRITE.getWidth(null) * PIXEL_SCALE / 2, y + item.SPRITE.getHeight(null) * PIXEL_SCALE + 32);

            x += item.SPRITE.getWidth(null) * PIXEL_SCALE + iPadding;
        }

        // Objekte
        super.draw(g);
    }
    
}
\end{lstlisting}

\subsection{LehrerObject}
Das LehrerObject ist ebenfalls eine Erweiterung des BaseObject und fügt die in der Lösungsidee besprochenen Eigenschaften von Schwierigkeit, Lebenspunkten, Fachschaft und ein Liste aus im Kampf zu verwendenen Items zu. Diese sind hier das Integer Level und hp (Healthpoints), der String Fachschaft und das Array Lehrer-Items. Diese werden dann zusammen mit den Parametern aus dem BaseObjekt im darauffolgenden Constructor initialisiert. Darauffolgend sind wieder die Funktionen für den Kampf.
\\
\begin{lstlisting}
    public class LehrerObject extends BaseObject{

    public final int LEVEL;
    public int hp;
    public final String FACHSCHAFT;
    public final List<ItemType> lehrer_items;

    public LehrerObject(Image sprite, int x, int y, String FACHSCHAFT, int hp, int LEVEL, List<ItemType> lehrer_items) {
        super(sprite, x, y);
        this.LEVEL = LEVEL;
        this.hp = hp;
        this.FACHSCHAFT = FACHSCHAFT;
        this.lehrer_items = lehrer_items;
    }

    public void verteidigung(int level) {
        // Es wird random ausgewählt, ob ein Item gewählt wird (und welches), oder ob der Lehrer nichts macht.
        Random rand = new Random();
        int move = rand.nextInt(5);
        int schaden;
        ItemType item_lehrer;

        if (move == 0) {
            schaden = level;
            item_lehrer = null;
        }
        else {
            item_lehrer = lehrer_items.get(move - 1);
            schaden = level - item_lehrer.LEVEL;
            if (schaden <= 0) {
                schaden = 0;
            }

        }

        hp -= schaden;

    }
    
    public void angriff() {
        Random rand = new Random();
        int move = rand.nextInt(4);
        int level;
        ItemType item_lehrer;
        
        item_lehrer = lehrer_items.get(move);
        level = item_lehrer.LEVEL;
        
        PlayerObject.INSTANCE.verteidigung(level);
    }

    @Override
    public void update(BaseScene scene) {
    }
}
\end{lstlisting}

\subsection{ItemType}
Die Items sind in der Form von im Glossar erklärten Enums gespeichert. Jedes Enum hat einen Namen und eine am Anfang festgelegte Anzahl von Parametern und deren Datentypen: der String NAME für den Namen des Items, das Image SPRITE für die Bilddatei und den Integer LEVEL für die Stärke des Items. 

\begin{lstlisting}
    public enum ItemType {
    ATOM("Atom", "/assets/items/atom.png", 0),
    BSOD("Error-Screen", "/assets/items/blue_screen_of_death.png", 2),
    BRENNER("Brenner", "/assets/items/brenner.png", 2),
    CAS("Taschenrechner", "/assets/items/cas.png", 2),
    CD("CD", "/assets/items/cd.png", 0),
    CHROME("Chrome", "/assets/items/chrome.png", 1),
    DNA("DNA", "/assets/items/dna.png", 1),
    FELDSTECHER("Feldstecher", "/assets/items/feldstecher.png", 1),
    GLUEHBIRNE("Glühbirne", "/assets/items/glühbirne.png", 2),
    IKOSAEDER("Ikosaeder", "/assets/items/ikosaeder.png", 1),
    KRAFTMESSER("Federkraftmesser", "/assets/items/kraftmesser.png", 0),
    LAUTSPRECHER("Lautsprecher", "/assets/items/lautsprecher.png", 0),
    LINEAL("Lineal", "/assets/items/lineal.png", 0),
    MAGNET("Magnet", "/assets/items/magnet.png", 1),
    MIKROSKOP("Mikroskop", "/assets/items/mikroskop.png", 2),
    NERV("Nervenzelle", "/assets/items/nerv.png", 0),
    NEWTONSAPFEL("Newtons Apfel", "/assets/items/newtons_apfel.png", 0),
    THALES("Satz des Thales", "/assets/items/satz_des_thales.png", 1),
    SCHUTZBRILLE("Schutzbrille", "/assets/items/schutzbrille.png", 1),
    SCHAEDEL("Schädel", "/assets/items/schädel.png", 1),
    SEIZUREDFROG("Sezierter Frosch", "/assets/items/seizured_frog.png", 2),
    SPRITZFLASCHE("Spritzflasche", "/assets/items/spritzflasche.png", 0),
    SAEURE("Säure", "/assets/items/säure.png", 2),
    TASTATUR("Tastatur", "/assets/items/tastatur.png", 2),
    THERMOMETER("Thermometer", "/assets/items/thermometer.png", 2),
    VIRUS("Virus", "/assets/items/virus.png", 0),
    WLAN("Wlan", "/assets/items/wlan.png", 1),
    WUNDERKERZE("Wunderkerze", "/assets/items/wunderkerze.png", 1),
    ZETTEL("Zettel", "/assets/items/zettel.png", 2),
    ZIRKEL("Zirkel", "/assets/items/zirkel.png", 0);
        
    public final String NAME;
    public final Image SPRITE;
    public final int LEVEL;
    
    ItemType(String name, String imagePath, int level) {
        this.NAME = name;
        this.LEVEL = level;
        this.SPRITE = Images.get(imagePath);
    }
\end{lstlisting}

Darauf folgt die Funktion, welche dem Spieler nach gewinnen eines Minispiels das Item zuweist.
Dies geschieht über zwei ineinander integrierte bereits erklärte switch-case-Selektionen, welche zuerst die Fachschaft des Minigames überprüfen und dann ein Item mit der dem Schwierigkeitsgrad des Minigames entsprechenden Stärke dem Inventar zuweisen.
\\
\begin{lstlisting}
      /**
     * Gibt ein Item als Belohnung für das Beenden eines Minispiels zurück.
     * @param minigame Das Minispiel, welches beendet wurde
     * @param level Das Level, welches beendet wurde
     */
    public static ItemType getMinigameReward(Department minigame, int level) {
        switch(minigame) {
            case COMPUTER_SCIENCE -> {
                switch(level) {
                    case 0 -> {
                        return ItemType.CD;
                    } case 1 -> {
                        return ItemType.CHROME;
                    } case 2 -> {
                        return ItemType.TASTATUR;
                    }
                }
            } case PHYSICS -> {
                switch(level) {
                    case 0 -> {
                        return ItemType.KRAFTMESSER;
                    } case 1 -> {
                        return ItemType.MAGNET;
                    } case 2 -> {
                        return ItemType.GLUEHBIRNE;
                    }
                }
            } case MATHEMATICS -> {
                switch(level) {
                    case 0 -> {
                        return ItemType.ZIRKEL;
                    } case 1 -> {
                        return ItemType.THALES;
                    } case 2 -> {
                        return ItemType.CAS;
                    }
                }
            } case BIOLOGY -> {
                switch(level) {
                    case 0 -> {
                        return ItemType.VIRUS;
                    } case 1 -> {
                        return ItemType.SCHAEDEL;
                    } case 2 -> {
                        return ItemType.SEIZUREDFROG;
                    }
                }
            } case CHEMISTRY -> {
                switch(level) {
                    case 0 -> {
                        return ItemType.SPRITZFLASCHE;
                    } case 1 -> {
                        return ItemType.SCHUTZBRILLE;
                    } case 2 -> {
                        return ItemType.SAEURE;
                    }
                }
            }  
        }
        return null;
        
    }
\end{lstlisting}

Der lezte Teil benutz eine Folge if-else-Abfragen, um den Lehrern im Kampf ihr Items zuzuweisen. Diese fragen zuerst nach der Fachschaft des Lehrers/ der Lehrerin und stellen diesen dann einen Pool aus Items bereit, welche sie dann im Kampf verwenden können.
\\
\begin{lstlisting}
     public static List<ItemType> getItems(int level, Department fachschaft) { // Items mit jeweiligen Leveln an Lehrer mit jeweiligen Leveln  verteilen
        if (fachschaft == Department.COMPUTER_SCIENCE) {
            if (level == 0) {
                return List.of(     //die Items werden zurückgegeben
                        ItemType.CD, //level 1
                        ItemType.LAUTSPRECHER, //level 1
                        ItemType.WLAN, //level 1 (eigentlich level 2)
                        ItemType.CHROME//level 2
                );
            } else if (level == 1) {
                return List.of(
                        ItemType.CD, //level 1
                        ItemType.LAUTSPRECHER, //level 1
                        ItemType.WLAN, //level 2
                        ItemType.CHROME //level 2
                );
            } else if (level == 2) {
                return List.of(
                        ItemType.LAUTSPRECHER, //level 1
                        ItemType.CHROME, //level 2
                        ItemType.WLAN, //level 2
                        ItemType.BSOD // level 3
                );

            }
        } else if (fachschaft == Department.PHYSICS) {
            if (level == 0) {
                return List.of(
                        ItemType.NEWTONSAPFEL, //level 1
                        ItemType.KRAFTMESSER, //level 1
                        ItemType.FELDSTECHER, //level 1 (eigentlich level 2)
                        ItemType.MAGNET//level 2
                );
            } else if (level == 1) {
                return List.of(
                        ItemType.NEWTONSAPFEL, //level 1
                        ItemType.KRAFTMESSER, //level 1
                        ItemType.FELDSTECHER, //level 2
                        ItemType.MAGNET //level 2
                );
            } else if (level == 2) {
                return List.of(
                        ItemType.NEWTONSAPFEL, //level 1
                        ItemType.FELDSTECHER, //level 2
                        ItemType.MAGNET, //level 2
                        ItemType.GLUEHBIRNE // level 3
                );

            }
        } else if (fachschaft == Department.MATHEMATICS) {
            if (level == 0) {
                return List.of(
                        ItemType.LINEAL, //level 1
                        ItemType.ZIRKEL, //level 1
                        ItemType.IKOSAEDER, //level 1 (eigentlich level 2)
                        ItemType.THALES//level 2
                );
            } else if (level == 1) {
                return List.of(
                        ItemType.LINEAL, //level 1
                        ItemType.ZIRKEL, //level 1
                        ItemType.IKOSAEDER, //level 2
                        ItemType.THALES //level 2
                );
            } else if (level == 2) {
                return List.of(
                        ItemType.LINEAL, //level 1
                        ItemType.IKOSAEDER, //level 2
                        ItemType.THALES, //level 2
                        ItemType.ZETTEL // level 3
                );

            }
        } else if (fachschaft == Department.BIOLOGY) {
            if (level == 0) {
                return List.of(
                        ItemType.VIRUS, //level 1
                        ItemType.NERV, //level 1
                        ItemType.SCHAEDEL, //level 1 (eigentlich level 2)
                        ItemType.DNA//level 2
                );
            } else if (level == 1) {
                return List.of(
                        ItemType.VIRUS, //level 1
                        ItemType.NERV, //level 1
                        ItemType.SCHAEDEL, //level 2
                        ItemType.DNA //level 2
                );
            } else if (level == 2) {
                return List.of(
                        ItemType.VIRUS, //level 1
                        ItemType.SCHAEDEL, //level 2
                        ItemType.DNA, //level 2
                        ItemType.SEIZUREDFROG // level 3
                );

            }
        } else if (fachschaft == Department.CHEMISTRY) {
            if (level == 0) {
                return List.of(
                        ItemType.SPRITZFLASCHE, //level 1
                        ItemType.ATOM, //level 1
                        ItemType.SCHUTZBRILLE, //level 1 (eigentlich level 2)
                        ItemType.WUNDERKERZE//level 2
                );
            } else if (level == 1) {
                return List.of(
                        ItemType.SPRITZFLASCHE, //level 1
                        ItemType.ATOM, //level 1
                        ItemType.SCHUTZBRILLE, //level 2
                        ItemType.WUNDERKERZE //level 2
                );
            } else if (level == 2) {
                return List.of(
                        ItemType.SPRITZFLASCHE, //level 1
                        ItemType.SCHUTZBRILLE, //level 2
                        ItemType.WUNDERKERZE, //level 2
                        ItemType.BRENNER // level 3
                );

            }
        }

        throw new IllegalArgumentException("Konnte nicht die Items für Fachschaft "+fachschaft+", Level "+level+" ermitteln!");
    }
}
\end{lstlisting}